\documentclass[11pt]{article}
\usepackage[a4paper,margin=11mm]{geometry}
\usepackage[no-math]{fontspec}
\setmainfont{Latin Modern Roman}
\usepackage{mathtools,amsmath,amssymb,amsfonts,amsthm,enumerate,arydshln,xurl}
\usepackage[unicode,hidelinks]{hyperref}
\usepackage{thmtools,thm-restate}
\setlength\emergencystretch{3em}
\declaretheorem[name=Theorem, parent=section]{theorem}
\declaretheorem[name=Corollary, sibling=theorem]{corollary}
\declaretheorem[name=Proposition, sibling=theorem]{proposition}
\declaretheorem[name=Lemma, sibling=theorem]{lemma}
\declaretheorem[name=Conjecture, sibling=theorem]{conjecture}
\declaretheorem[name=Observation, sibling=theorem]{observation}
\declaretheorem[name=Fact, sibling=theorem]{fact}
\declaretheorem[name=Definition, sibling=theorem, style=definition]{definition}
\declaretheorem[name=Remark, sibling=theorem, style=definition]{remark}
\declaretheorem[name=Problem, sibling=theorem, style=definition]{problem}
\declaretheorem[name=Example, sibling=theorem, style=definition]{example}
\declaretheorem[name=Claim, sibling=theorem, style=remark]{claim}

\DeclareMathOperator{\Diag}{Diag}
\DeclareMathOperator{\supp}{supp}
\DeclareMathOperator{\subrank}{Q}
\DeclareMathOperator{\slicerank}{SR}
\DeclareMathOperator{\bordersubrank}{\underline{Q}}
\DeclareMathOperator{\tensorrank}{R}
\newcommand{\FF}{\mathbb{F}}
\newcommand{\KK}{\mathbb{K}}
\newcommand{\RR}{\mathbb{R}}
\newcommand{\CC}{\mathbb{C}}
\newcommand{\NN}{\mathbb{N}}
\newcommand{\ZZ}{\mathbb{Z}}
\newcommand{\eps}{\varepsilon}
\DeclareMathOperator{\GL}{GL}

\DeclareMathOperator{\Mat}{Mat}

\DeclareMathOperator{\crk}{crk}
\DeclareMathOperator{\ncrk}{ncrk}
\DeclareMathOperator{\linspan}{span}
\DeclareMathOperator{\rank}{rank} \DeclareMathOperator{\col}{colspan} 

%\DeclareMathOperator{\Diag}{Diag}
%

\DeclareMathOperator{\maxrank}{max-rank}
\DeclareMathOperator{\minrank}{min-rank}
\DeclareMathOperator{\minrnk}{min-rank}
\DeclareMathOperator{\mincov}{min-cov}
\DeclareMathOperator{\minsupp}{min-supp}
\DeclareMathOperator{\maxsupp}{max-supp}


\newcommand{\id}{\mathrm{Id}}
\newcommand{\overCl}[1]{{{#1}^{\overline{\FF}}}}
\newcommand{\overK}[1]{{#1^{\mathbb{K}}}}

% \makeatletter
% \newcommand{\longdash}[1][2em]{%
%   \makebox[#1]{$\m@th\smash-\mkern-7mu\cleaders\hbox{$\mkern-2mu\smash-\mkern-2mu$}\hfill\mkern-7mu\smash-$}}
% \makeatother

%
  \newcommand{\beq}{\begin{equation}}
  \newcommand{\eeq}{\end{equation}}
  \newcommand{\beqn}{\begin{equation*}}
  \newcommand{\eeqn}{\end{equation*}}
  \newcommand{\beqr}{\begin{eqnarray}}
  \newcommand{\eeqr}{\end{eqnarray}}
  \newcommand{\beqrn}{\begin{eqnarray*}}
  \newcommand{\eeqrn}{\end{eqnarray*}}

\IfFontExistsTF{NimbusRoman-Regular}{\newfontface\FixedOriginalUprightFace{NimbusRoman-Regular}}{\PackageError{fixed-source-operators}{Required named NimbusRoman-Regular face missing}{No fallback permitted}}
\makeatletter
% Three finite zero-argument glyphs only. No underaccent or wtilde API is implemented.
\edef\FixedOriginalLowerTilde{\noexpand\mathchar\number\numexpr\symlargesymbols*256+101\relax}
\newcommand{\FixedQGlyphStyle}[2]{\begingroup\setbox0=\hbox{$\m@th#1\text{\normalfont\FixedOriginalUprightFace Q}$}\setbox2=\hbox{$\m@th#1\FixedOriginalLowerTilde$}\dimen0=\dp0\advance\dimen0 by\ht2\advance\dimen0 by1.23pt\typeout{FIXEDGLYPH Q STYLE\meaning#1 BASEwidth=\the\wd0 BASEheight=\the\ht0 BASEdepth=\the\dp0 TILDEwidth=\the\wd2 TILDEheight=\the\ht2 TILDEdepth=\the\dp2 LOWER=\the\dimen0}\hbox{\rlap{\hbox to\wd0{\hss\lower\dimen0\copy2\hss}}\box0}\endgroup}
\newcommand{\FixedRGlyphStyle}[2]{\begingroup\setbox0=\hbox{$\m@th#1\text{\normalfont\FixedOriginalUprightFace R}$}\setbox2=\hbox{$\m@th#1\FixedOriginalLowerTilde$}\dimen0=\dp0\advance\dimen0 by\ht2\advance\dimen0 by1.23pt\typeout{FIXEDGLYPH R STYLE\meaning#1 BASEwidth=\the\wd0 BASEheight=\the\ht0 BASEdepth=\the\dp0 TILDEwidth=\the\wd2 TILDEheight=\the\ht2 TILDEdepth=\the\dp2 LOWER=\the\dimen0}\hbox{\rlap{\hbox to\wd0{\hss\lower\dimen0\copy2\hss}}\box0}\endgroup}
\newcommand{\FixedSRGlyphStyle}[2]{\begingroup\setbox0=\hbox{$\m@th#1\text{\normalfont\FixedOriginalUprightFace SR}$}\setbox2=\hbox{$\m@th#1\FixedOriginalLowerTilde$}\dimen0=\dp0\advance\dimen0 by\ht2\advance\dimen0 by1.23pt\typeout{FIXEDGLYPH SR STYLE\meaning#1 BASEwidth=\the\wd0 BASEheight=\the\ht0 BASEdepth=\the\dp0 TILDEwidth=\the\wd2 TILDEheight=\the\ht2 TILDEdepth=\the\dp2 LOWER=\the\dimen0}\hbox{\rlap{\hbox to\wd0{\hss\lower\dimen0\copy2\hss}}\box0}\endgroup}
\DeclareMathOperator{\asympsubrank}{\mathpalette\FixedQGlyphStyle\relax}
\DeclareMathOperator{\asymprank}{\mathpalette\FixedRGlyphStyle\relax}
\DeclareMathOperator{\asympslicerank}{\mathpalette\FixedSRGlyphStyle\relax}
\makeatother
\begin{document}\begin{center}\Large Problem 2472: Cube-root asymptotic subrank lower bound\end{center}\noindent\textbf{Authors:} Jop Briët, Matthias Christandl, Itai Leigh, Amir Shpilka and Jeroen Zuiddam\par\noindent\textbf{Work/version:} Discreteness of Asymptotic Tensor Ranks; exact arXiv2306.01718v4 (9 September 2025); native CC BY4.0; matching Discrete Analysis2025:15, 51pp. (10 September 2025), DOI10.19086/da.143834; printedPDFCC BY3.0\par\noindent\textbf{Native source/license:} \url{https://arxiv.org/src/2306.01718v4}, \url{https://creativecommons.org/licenses/by/4.0/}\par\noindent\textbf{Difficulty:} H2: The argument relates maximal ranks in different slice directions through simultaneous polynomially generic column selections. Conciseness forces the product of each pair of slice-span ranks to dominate the remaining dimension. A third tensor power restricts to a matrix-multiplication tensor whose asymptotic subrank transfers those pairwise bounds back to the original tensor.\par\medskip\noindent\textbf{Editorial source boundary:} Theorem 2.11, identical to announced 1.8, is the sole selected cube-root outcome. The complete max-rank column-selection/polynomial/conciseness argument, exact matrix multiplication and tensor restriction definitions, third-power restriction and full all-field closing chain are retained. The fourth-root result is necessary author field-extension framing and uncounted context. Named Schwartz-Zippel, Strassen field-extension invariance and matrix-multiplication asymptotic subrank remain precise established prerequisites. Exact native v4 CC BY 4.0 and published PDF CC BY 3.0 identities remain separate. Original accents package acquisition is held after HTTP402; no package is cloned. Only the exact fixed zero-argument Q/R/SR under-tilde operators have an approved independent typography resolver using the original installed CMEX0x65 glyph and explicitly named installed NimbusRoman-Regular upright bases. Literal author calls/formulas and source helper definitions remain byte-bound; no arbitrary accent API, unknown argument or font fallback is supported. No mathematical proof, hypothesis or parameter is changed or generated. Upstream classes/logos/archive stay private and unexecuted. Full packet qualification remains pending.\par\medskip\hrule\medskip\noindent\textbf{Original human-authored mathematics follows.}\par
% Exact native author source LF lines 885--978
\setcounter{section}{1}\setcounter{subsection}{3}\setcounter{theorem}{14}
\subsection{Tensor preliminaries}\label{subsec:prelim}

We summarize in this section some standard tensor preliminaries and terminology that we will use throughout the paper, following the papers of Strassen \cite{MR882307,strassen1988asymptotic,strassen1991degeneration} and the book of Bürgisser, Clausen and Shokrollahi~\cite{burgisser1997algebraic}.


\subsubsection*{Tensors}
%
%
Let $\FF$ be any field. For any $n \in \NN$ we denote by $e_1, \ldots, e_n$ the standard basis elements of the vector space $\FF^n$. For any $n_1, n_2, n_3 \in \NN$, we denote by $\FF^{n_1} \otimes \FF^{n_2} \otimes \FF^{n_3}$ the vector space of order-three tensors $T = \sum_{i \in [n_1]} \sum_{j \in [n_2]} \sum_{k \in [n_3]} T_{i,j,k}\, e_i \otimes e_j \otimes e_k$ with coefficients $T_{i,j,k} \in \FF$. In most of this paper we will be dealing only with order-three tensors, so instead of order-three tensor we will just say tensor. 

%
\subsubsection*{Restriction and equivalence} 
We will be using a natural preordering on tensors called restriction, which is defined as follows. For any two tensors $T \in \FF^{n_1} \otimes \FF^{n_2} \otimes \FF^{n_3}$ and $S \in \FF^{m_1} \otimes \FF^{m_2} \otimes \FF^{m_3}$, we say $T$ \emph{restricts to} $S$ and write $T \geq S$ if there are linear maps $L_\ell : \FF^{n_\ell} \to \FF^{m_\ell}$ such that $(L_1 \otimes L_2 \otimes L_3)T = S$. We say $S$ and $T$ are \emph{equivalent} if $T \geq S$ and $S \geq T$. Most tensor properties we will study are invariant under equivalence and monotone under restriction (see \autoref{lem:monotone-restrict}). Finally, we say $S, T \in \FF^{n_1} \otimes \FF^{n_2} \otimes \FF^{n_3}$ are \emph{isomorphic} if $(L_1 \otimes L_2 \otimes L_3)T = S$ for invertible linear maps $L_\ell : \FF^{n_\ell} \to \FF^{n_\ell}$.
%

\subsubsection*{Flattenings and flattening ranks}
Let $T \in V_1 \otimes V_2 \otimes V_3$. Grouping together $V_2$ and $V_3$ gives an element  $T_{1} \in V_1 \otimes (V_2 \otimes V_3)$, and we similarly obtain $T_2$ and $T_3$. %(for any choice of distinct ${\ell_1},{\ell_2},{\ell_3} \in [3]$). 
%We may think of $T_{\ell_1}$ as an $n_{\ell_1} \times (n_{\ell_2}\cdot n_{\ell_3})$ matrix. We call the $T_{\ell_1}$ 
We may think of the $T_\ell$ as matrices and we call them 
the \emph{flattenings} of $T$.
Let $\tensorrank_{\ell}(T)$ ve the matrix rank of $T_{\ell}$. We call $\tensorrank_1(T)$, $\tensorrank_2(T)$ and $\tensorrank_3(T)$ 
%\refnote{7}{missing plural " 's"}\ILnote{added it}\Jshort{I don't like the plural s so much after symbols so I solved it differently} 
the \emph{flattening ranks} of $T$.
%Equivalently, we may slice $T$ into matrices $A_i = (T_{i,j,k})_{j,k} \in \FF^{n_2} \otimes \FF^{n_3}$, $B_j = (T_{i,j,k})_{i,k} \in \FF^{n_1} \otimes \FF^{n_3}$ and $C_k = (T_{i,j,k})_{i,j} \in \FF^{n_1} \otimes \FF^{n_2}$
The triple $(\tensorrank_1(T), \tensorrank_2(T), \tensorrank_3(T))$ is often called the \emph{multilinear rank} of $T$. 


\subsubsection*{Conciseness} 
%\refnote{3}{concise actual definition, out of: \phantomsection\label{ref:3-concise-def13}, \ref{ref:3-first-concise-use},\ref{ref:3-first-concise-def},\ref{ref:3-concise-use7},\ref{ref:3-concise-def8},\ref{ref:3-concise-use10},\ref{ref:3-concise-def13}}
Let $T \in \FF^{n_1} \otimes \FF^{n_2} \otimes \FF^{n_3}$. From the definition of flattening rank it follows that $\tensorrank_\ell(T) \leq n_\ell$ for every $\ell \in [3]$. We call $T$ \emph{concise} if $\tensorrank_\ell(T) = n_\ell$  for every~$\ell \in [3]$. Conciseness essentially says that the tensor cannot fit in a smaller tensor space. The property of being concise is very mild in the following sense. We call a tensor $S$ a \emph{subtensor} of $T = \sum_{i,j,k} T_{i,j,k}\, e_i \otimes e_j \otimes e_k$ if $S$ is of the form $S = \sum_{i\in I,j\in J,k\in K} T_{i,j,k}\, e_i \otimes e_j \otimes e_k$ for some subsets $I,J,K$.

\begin{lemma}\label{lem:equiv-concise}
    Every tensor is equivalent to a concise tensor, which is unique up to equivalence. Concretely, every tensor is equivalent to a subtensor that is concise.
\end{lemma}
\begin{proof}
    The first statement is a standard fact \cite[Section 14.3]{burgisser1997algebraic}. For the second statement, let $T = \sum_{i,j,k} T_{i,j,k}\, e_i \otimes e_j \otimes e_k \in \FF^{n_1} \otimes \FF^{n_2} \otimes \FF^{n_3}$. Let $A_i = (T_{i,j,k})_{j,k} \in \FF^{n_2} \otimes \FF^{n_3}$ for $i \in [n_1]$ be the 1-slices of~$T$. Then one verifies that the dimension of the span of the~$A_i$ equals the flattening rank $r_1=\tensorrank_1(T)$. Suppose $A_1, \ldots, A_r$ are linearly independent, and let $S$ be the tensor with these matrices as 1-slices. Then $S$ is a subtensor of $T$ so clearly $T \geq S$. Also, since the matrices $A_{r+1}, \ldots, A_{n_1}$ are in the span of the matrices $A_1, \ldots, A_r$ we have that $S \geq T$. Since $S$ and $T$ are thus equivalent, their flattening ranks are equal. Repeating this construction for the other two directions gives the claimed subtensor.
    %
    %Zij T een n1 x n2 x n3 tensor. Zij A_1, ..., A_{n1} de 1-slices van T, dus T = \sum_{i=1}^{n1} e_i \otimes A_i. De flattening rank R1(T) is de dimensie van de span van de A_i. Stel A_1 is in de span van A_2, ..., A_{n1}. Laat S = \sum_{i=2}^{n1} e_i \otimes A_i. Dit is een subtensor van T. Er geldt T \geq S en S \geq T. De eerste restrictie T \geq S is duidelijk, want S is een subtensor van T. De tweede restrictie is makkelijk in te zien de slices van S zijn A_2, ..., A_{n1} en A_1 is in de span van die matrices. Met andere woorden S en T zijn equivalent, en dus zijn hun flattening ranks hetzelfde (want flattening ranks zijn monotoon onder restrictie). Met deze procedure kunnen we een tensor U vinden die een subtensor is van T, equivalent is aan T, en concise is.
\end{proof}

Not every space $\FF^{n_1} \otimes \FF^{n_2} \otimes \FF^{n_3}$ contains concise tensors, as can be seen from the following implication (which is in fact known to be an equivalence \cite[Theorem~2]{MR2776439}):

\begin{lemma}\label{cl:concise-formats}
For any $n_1,n_2,n_3\in\NN$, if there is a concise tensor in $\FF^{n_1} \otimes \FF^{n_2} \otimes \FF^{n_3}$, then $n_1 \leq n_2\cdot n_3$, $n_2 \leq n_1\cdot n_3$, and $n_3 \leq n_1\cdot n_2$.
\end{lemma}
\begin{proof}
    Suppose $T \in \FF^{n_1} \otimes \FF^{n_2} \otimes \FF^{n_3}$ is concise. Then $n_{\ell_1} = \rank(T_{\ell_1})$ where $T_{\ell_1}$ is the previously defined flattening of $T$. Since $T_{\ell_1}$ is an $n_{\ell_1} \times (n_{\ell_2}\cdot n_{\ell_3})$ matrix, we have $\tensorrank(T_{\ell_1}) \leq n_{\ell_2} \cdot n_{\ell_3}$. Thus $n_{\ell_1} \leq n_{\ell_2} \cdot n_{\ell_3}$ for every choice of distinct ${\ell_1},{\ell_2},{\ell_3} \in [3]$.
\end{proof}
%
%
%
%

%
%
%
%
%
%
%
%

\subsubsection*{Unit tensors, tensor rank, subrank, and slice rank}
For every $r \in \NN$ we define the tensor $\langle r \rangle = \sum_{i=1}^r e_i \otimes e_i \otimes e_i \in \FF^r \otimes \FF^r \otimes \FF^r$. We call~$\langle r \rangle$ the \emph{unit tensor} of rank $r$.
We say a tensor has \emph{tensor rank one} if it is of the form $u \otimes v \otimes w$ for nonzero vectors $u,v,w$.
The tensor rank of a tensor $T$ is the smallest number $r$ such that $T$ can be written as a sum of $r$ tensors with tensor rank one. 
%
We denote tensor rank by $\tensorrank(T)$. Equivalently, the tensor rank $\tensorrank(T)$ is the smallest number $r$ such that $T \leq \langle r\rangle$. The \emph{subrank} of $T$ is the largest number $r$ such that $\langle r\rangle \leq T$. We denote subrank by $\subrank(T)$. 
%\ILnote{maybe we should add: "For $\ell\in[3]$, we denote by $\subrank_\ell(T)$ the maximal rank of any linear combination of the $\ell$-slices of $T$ (see \autoref{subsec:maxrank}". without it -- we haven't defined it yet but it appears in the lemma below.} 
We say a tensor has \emph{slice rank one} if any of its flattening ranks equals~one. The slice rank of a tensor $T$ is the smallest number $r$ such that $T$ can be written as a sum of $r$ tensors with slice rank one.
%
Subrank, slice rank, the flattening ranks and tensor rank satisfy $\subrank(T) \leq \slicerank(T) \leq \tensorrank_\ell(T) \leq \tensorrank(T)$ for every $\ell\in [3]$.

It is also easy to see that these measures are monotone under restrictions:

\begin{lemma}\label{lem:monotone-restrict}
    For any two tensors $T$ and $S$, if $T\geq S$, then we have $\tensorrank(T)\geq \tensorrank(S)$, $\slicerank(T)\geq \slicerank(S)$, $\subrank(T)\geq \subrank(S)$, and for every $\ell\in[3]$, %$\subrank_\ell(T)\geq \subrank_\ell(S)$ and 
    $\tensorrank_\ell(T)\geq \tensorrank_\ell(S)$.
\end{lemma}


\subsubsection*{Kronecker product and asymptotic tensor parameters}
%
Asymptotic tensor parameters are defined by ``regularizing'' a tensor parameter over large powers under the Kronecker product. For any two tensors 
\begin{align*}
T &= \sum_{i,j,k} T_{i,j,k}\, e_i \otimes e_j \otimes e_k \in \FF^{n_1} \otimes \FF^{n_2} \otimes \FF^{n_3}\\
S &= \sum_{a,b,c} S_{a,b,c} \, e_a\otimes e_b \otimes e_c \in \FF^{m_1} \otimes \FF^{m_2} \otimes \FF^{m_3}
\end{align*}
the Kronecker product $T \boxtimes S \in \FF^{n_1 m_1} \otimes \FF^{n_2 m_2} \otimes \FF^{n_3 m_3}$ is defined as
\[
T \boxtimes S = \sum_{i,j,k} \sum_{a,b,c} T_{i,j,k} S_{a,b,c}\, (e_i \otimes e_{a}) \otimes (e_j \otimes e_{b}) \otimes (e_k \otimes e_{c}).
\]
We define $T^{\boxtimes n}$ as the product of $n$ copies of $T$. The asymptotic tensor rank is defined as $\asymprank(T) = \lim_{n \to \infty} \tensorrank(T^{\boxtimes n})^{1/n}$, which, because of sub-multiplicativity of $\tensorrank$, equals the infimum $\inf_{n} \tensorrank(T^{\boxtimes n})^{1/n}$ (by Fekete's lemma). The asymptotic subrank is defined as $\asympsubrank(T) = \lim_{n \to \infty} \subrank(T^{\boxtimes n})^{1/n}$, which, because of super-multiplicativity of the subrank $\subrank$, equals the supremum $\sup_{n} \subrank(T^{\boxtimes n})^{1/n}$. Over the complex numbers, the asymptotic slice rank is defined as $\asympslicerank(T) = \lim_{n\to\infty} \slicerank(T^{\boxtimes n})^{1/n}$. By the characterization in \cite{MR4495838} this limit exists. Over other fields, we generally do not know whether the limit $\lim_{n\to\infty} \slicerank(T^{\boxtimes n})^{1/n}$ exists, except for oblique tensors. When we do not know whether the limit exists we will instead consider the liminf and refer to this as the ``asymptotic slice rank''.\footnote{The results we have over finite fields for asymptotic slice rank hold equally well using liminf or limsup.}

% Exact native author source LF lines 1001--1001
\setcounter{section}{1}
\section{Max-rank of slice spans and asymptotic subrank lower bound}\label{sec:max-rank}
% Exact native author source LF lines 1019--1209

\subsection{Max-rank of matrix subspaces and slice spans of a tensor}\label{subsec:maxrank}

%
    %
    We denote by $\FF^{n_1} \otimes \FF^{n_2}$ the space of $n_1 \times n_2$ matrices over $\FF$.
    We define the \emph{max-rank} of any matrix subspace $\mathcal{A} \subseteq \FF^{n_1} \otimes \FF^{n_2}$ as
    %
    \[
\maxrank(\mathcal{A}):=\max\{\rank(A):A\in\mathcal{A}\}.
    \]
    %
    %
    %
    %
%
%
For any tensor $T = \sum_{i,j,k} T_{i,j,k}\, e_i \otimes e_j \otimes e_k \in \FF^{n_1} \otimes \FF^{n_2} \otimes \FF^{n_3}$ we define the matrices %
%
%
%
%
%
%
\begin{align*}
T^{(1)}_i &= \sum_{j,k} T_{i,j,k}\, e_j \otimes e_k \quad (i \in [n_1])\\
T^{(2)}_j &= \sum_{i,k} T_{i,j,k}\, e_i \otimes e_k \quad (j \in [n_2])\\
T^{(3)}_k &= \sum_{i,j} T_{i,j,k}\, e_i \otimes e_j \quad (k \in [n_3]).
\end{align*}
For every $\ell \in [3]$, we call $\{T^{(\ell)}_i : i \in [n_\ell] \}$ the \emph{$\ell$-slices} of $T$. %
%
%
%
%
%
%
%
%
%
%
%
%
%
%
%
%
%
%
%
%
%
%
%
%
%
%
%
%
%
%
%
%
%
%
%
%
%
%
%
%
%
%
We define the $\ell$-slice span of $T$ as the matrix subspace
$\mathcal{A}_\ell = \linspan \{T^{(\ell)}_i : i \in [n_\ell] \}$
and denote its max-rank by
\[
\subrank_\ell(T) = \maxrank(\mathcal{A}_\ell).
\]
%
%
We note that the flattening rank $\tensorrank_\ell(T)$ defined earlier, is equal to the dimension of the subspace $\mathcal{A}_\ell$. 

%



%
%



%
%

%
%
%
%
%
%
%
%

%
%
%

%
%
%
%
%
%
%
%
%
%
%
%
%
%
%
%

%

%
%
%
%
%
%
%
%
%
%
%
%
%




\subsection{Lower bound on products of max-ranks of slice spans}\label{subsec:QiQj}

%

%

%

%

%\refnote{9}{The following paragraph sounds like a repetitions of earlier things. also \phantomsection\label{ref:9-repetition-22}\ref{ref:9-repetition-21}}
In this section we prove the aforementioned fundamental property of the max-ranks of the slice spans of a tensor, showing that they cannot all be small. This will be a crucial ingredient for proving the asymptotic subrank lower bound afterwards.

\begin{theorem}\label{th:uncertainty-principle-n}
    Let $T \in \FF^{n_1} \otimes \FF^{n_2} \otimes \FF^{n_3}$ be concise. Let ${\ell_1}, {\ell_2}, {\ell_3} \in [3]$ be distinct. Suppose that $|\FF|>n_{\ell_3}$. Then 
    \[
    \subrank_{\ell_1}(T)\subrank_{\ell_2}(T) \geq n_{\ell_3}.
    \]
\end{theorem}

%

%

%

A corollary of this is that there are two distinct directions with large max-rank:

\begin{corollary} %
\label{cor:twoLargeRankDirs}
    Let $T\in\FF^{n_1}\otimes\FF^{n_2}\otimes\FF^{n_3}$ be concise.
    Suppose $n_1 \leq n_2 \leq n_3$. Suppose that $|\FF| > n_3$. 
    %
    Then there are distinct ${\ell_1},{\ell_2}\in[3]$ such that
    \begin{align*}\subrank_{\ell_1}(T) &\geq \sqrt{n_3},\\
    \subrank_{\ell_2}(T) &\geq \sqrt{n_1}.
    \end{align*}
\end{corollary}
%
\begin{proof}%
    %
    We have by \autoref{th:uncertainty-principle-n} that
    \begin{align*}
    \subrank_1(T) \subrank_2(T) &\geq n_3\\
    \subrank_2(T) \subrank_3(T) &\geq n_1\\
    \subrank_1(T) \subrank_3(T) &\geq n_2 \geq n_1.
    \end{align*}
    Thus one of $\subrank_1(T)$ and $\subrank_2(T)$ is at least $\sqrt{n_3}$, and either $\subrank_3(T)$ is at least~$\sqrt{n_1}$ or both $\subrank_1(T)$ and $\subrank_2(T)$ are at least $\sqrt{n_1}$.
\end{proof}
% Exact native author source LF lines 1310--1471

For matrices $A \in \FF^{n\times m}$ and $B \in \FF^{n \times k}$ we denote by $[A; B] \in \FF^{n \times (m+k)}$ the concatenation of $A$ and $B$. For any matrix $A$ we denote by $A|_{[d]}$ the matrix consisting of the first $d$ columns of $A$. We denote by $\col(A)$ the span of the columns of $A$.

%\ILnote{Comment 9. Fixed below the instances of $[A;(BU)|_{[s]}$ which all missed the closing square bracket (missed it, and not placed the restriction inside, as the comment says).}
\begin{lemma}\label{lem:colspan}
    Let $A \in \FF^{n\times m}$ and $B \in \FF^{n \times k}$. There is a nonzero %$k^2$-variable 
    polynomial $p\in \FF[x_{1,1},x_{1,2},,\ldots,x_{i,j},\dots,x_{k,k}]$ of degree $s = \rank([A;B]) - \rank(A)$ such that for any matrix $U \in \FF^{k\times k}$ satisfying $p(U)\neq 0$, we have that
    \[
    \col([A;B]) = \col([A; (BU)|_{[s]}]).
    \]
\end{lemma}
\begin{proof}
    Let $r = \rank([A;B])$. Then, there is an $r\times r$ submatrix of $[A;B]$ with rank~$r$ induced by $r-s$ columns of~$A$ and~$s$ columns of~$B$.
    Let~$R\subseteq [n]$ be the set of row indices of this submatrix, and let $C_A\subseteq[m], C_B\subseteq \{m+1,\dots,m+k\}$ be the respective column indices of~$A$ and~$B$ in this submatrix.
    Let $p(X)$ be the determinant of the submatrix of $[A;BX]$ induced by the rows in~$R$ and columns in $C_A\cup \{m+1,\dots,m+s\}$, where $X = (x_{i,j})_{i,j=1}^k$ is a matrix of variables.
    
    We claim that~$p$ is nonzero. There is a permutation matrix~$V\in \FF^{k \times k}$ that swaps the columns in~$C_B$ with the columns in $\{m+1,\dots, m+s\}$. Since the submatrix of $[A;BV]$ induced by the rows in~$R$ and columns in $C_A\cup \{m+1,\dots,m+s\}$ is then of full rank, it follows that~$p(V)\ne 0$. 
        
    Let $U \in \FF^{k \times k}$ be a matrix such that $p(U) \neq 0$. 
    Then we have $\rank([A; (BU)|_{[s]}]) \geq r$. From $\col([A; (BU)|_{[s]}])\subseteq \col([A;B])$ and $\rank([A;B]) = r$, we find the equality $\col([A;B]) = \col([A; (BU)|_{[s]}])$.
\end{proof}

% \refnote{10}{Wouldn't it be more intuitive to merge 2.4-2.5 into one lemma where you directly add $s_i$ columns at a time?}

% \JBnote{The proofs were merged and only the second lemma remains, as suggested.}

% \begin{lemma}\label{lem:schwartz-zippel}
%     Let $A_1, \ldots, A_\ell \in \FF^{n \times k}$ be matrices. Suppose $|\FF| > \rank([A_1; \cdots; A_\ell])$. Let 
%     \[
%     s_i = \rank([A_1; \cdots ; A_{i}]) - \rank([A_1; \cdots ; A_{i-1}]).
%     \]
%     There is a matrix $U \in \FF^{k \times k}$ such that for all $i \in [\ell]$ we have 
%     \[
%     \col([A_1 ; A_2 ; \cdots ; A_{i-1} ; (A_i U)|_{[s_i]}]) = \col([A_1; \cdots ; A_i]).
%     \]
% \end{lemma}
% \begin{proof}
% From Lemma~\ref{lem:colspan}, we get that for each $i\in [\ell]$, there exists a nonzero~$k^2$-variable polynomial $p_i(X)\in \FF[x_{1,1},x_{1,2},\dots,x_{k,k}]$ of degree~$s_i$ such that for any $U\in \FF^{k\times k}$ satisfying~$p_i(U)\ne 0$,
% \[
%     \col([A_1 ; A_2 ; \cdots ; A_{i-1} ; (A_i U)|_{[s_i]}]) = \col([A_1; \cdots ; A_i]).
%     \]
% Let $q(X)=  p_1(X)\cdots p_\ell(X)$.
% This is a nonzero polynomial of degree $s_1 + \cdots +s_\ell = \rank([A_1; \cdots; A_\ell])<|\FF|$.
% It follows from the Schwartz--Zippel lemma \cite{DEMILLO1978193, schwartz1980fast, zippel1979probabilistic} that there is a~$U\in \FF^{k\times k}$ such that~$q(U) \ne 0$.
% This is to say that for all~$i\in[\ell]$ simultaneously,~$p_i(U)\ne 0$, which implies the result.
% \end{proof}

\begin{lemma}\label{lem:max-rank-induction}
    Let $A_1, \ldots, A_m \in \FF^{n \times k}$ be matrices. Suppose $|\FF| > \rank([A_1; \cdots; A_m])$. Let 
    \[
    s_i = \rank([A_1; \cdots ; A_{i}]) - \rank([A_1; \cdots ; A_{i-1}]).
    \]
    There is a matrix $U \in \FF^{k \times k}$ such that
\begin{equation}\label{eq:cspan2}
    \col([(A_1 U)|_{[s_1]} ; (A_2 U)|_{[s_2]} ; \cdots ; (A_m U)|_{[s_m]}]) = \col([A_1; \cdots ; A_m]).
\end{equation}
\end{lemma}
% \begin{proof}
% %
% Let~$U\in \FF^{k\times k}$ be a matrix as in Lemma~\ref{lem:schwartz-zippel}. We give a proof by induction.
% The base case is $\col((A_1 U)|_{[s_1]}) = \col(A_1)$,
% which follows directly from \autoref{lem:schwartz-zippel}.
% For the induction step, we assume
% \begin{equation}\label{eq:ih}
%     \col([A_1 ; A_2 ; \cdots ; A_i]) = \col([(A_1U)|_{[s_1]} ; \cdots ; (A_iU)|_{[s_i]}]).
% \end{equation}
% We now use that if $\col(A) = \col(B)$, then $\col([A ; C]) = \col([B;C])$ holds for any matrices $A,B,C$ with the same number of rows. Applying this to \eqref{eq:ih} gives
% \begin{multline}\label{eq:x}
%     \col([A_1 ; A_2 ; \cdots ; A_i; (A_{i+1} U)|_{[s_{i+1}]}])\\ = \col([(A_1U)|_{[s_1]} ; \cdots ; (A_iU)|_{[s_i]}; (A_{i+1} U)|_{[s_{i+1}]}]).
% \end{multline}
% From \autoref{lem:schwartz-zippel} we know that 
% \begin{equation}\label{eq:sz}
%     \col([A_1; \cdots ; A_{i+1}]) = \col([A_1 ; A_2 ; \cdots ; A_i ; (A_{i+1} U)|_{[s_{i+1}]}]).
% \end{equation}
% %
% %
% Then \eqref{eq:x} and \eqref{eq:sz} together give
%     \[
%     \col([A_1 ; A_2 ; \cdots ; A_i; A_{i+1}]) = \col([(A_1U)|_{[s_1]} ; \cdots ; (A_iU)|_{[s_i]}; (A_{i+1}U)|_{[s_{i+1}]}]).
%     \]
% This proves the claim by induction.
% \end{proof}


\begin{proof}
We begin by showing that there exists a matrix $U \in \FF^{k \times k}$ such that for all $i \in [m]$ we have 
\begin{equation}\label{eq:cspan1}
    \col([A_1 ; A_2 ; \cdots ; A_{i-1} ; (A_i U)|_{[s_i]}]) = \col([A_1; \cdots ; A_i]).
\end{equation}
From Lemma~\ref{lem:colspan}, we get that for each $i\in [m]$, there exists a nonzero~$k^2$-variable polynomial $p_i(X)\in \FF[x_{1,1},x_{1,2},\dots,x_{k,k}]$ of degree~$s_i$ such that for any $U\in \FF^{k\times k}$ satisfying~$p_i(U)\ne 0$,
\[
    \col([A_1 ; A_2 ; \cdots ; A_{i-1} ; (A_i U)|_{[s_i]}]) = \col([A_1; \cdots ; A_i]).
    \]
Let $q(X)=  p_1(X)\cdots p_m(X)$.
This is a nonzero polynomial of degree $s_1 + \cdots +s_m = \rank([A_1; \cdots; A_m])<|\FF|$.
It follows from the Schwartz--Zippel lemma \cite{DEMILLO1978193, schwartz1980fast, zippel1979probabilistic} that there is a~$U\in \FF^{k\times k}$ such that~$q(U) \ne 0$.
This is to say that for all~$i\in[m]$ simultaneously,~$p_i(U)\ne 0$, which implies the existence of the desired matrix~$U$.

Let~$U\in \FF^{k\times k}$ be a matrix so that~\eqref{eq:cspan1} holds for each $i\in [m]$.
Next, we show by induction on~$i\in[m]$ that~$U$ also satisfies~\eqref{eq:cspan2}.
%Let~$U\in \FF^{k\times k}$ be a matrix as in Lemma~\ref{lem:schwartz-zippel}. We give a proof by induction.
The base case is $\col((A_1 U)|_{[s_1]}) = \col(A_1)$,
which follows directly.
For the induction step, we assume
\begin{equation}\label{eq:ih}
    \col([A_1 ; A_2 ; \cdots ; A_i]) = \col([(A_1U)|_{[s_1]} ; \cdots ; (A_iU)|_{[s_i]}]).
\end{equation}
We now use that if $\col(A) = \col(B)$, then $\col([A ; C]) = \col([B;C])$ holds for any matrices $A,B,C$ with the same number of rows. Applying this to \eqref{eq:ih} gives
\begin{equation}\label{eq:x}
    \col([A_1 ; A_2 ; \cdots ; A_i; (A_{i+1} U)|_{[s_{i+1}]}]) = \col([(A_1U)|_{[s_1]} ; \cdots ; (A_iU)|_{[s_i]}; (A_{i+1} U)|_{[s_{i+1}]}]).
\end{equation}
From \eqref{eq:cspan1} we know that 
\begin{equation}\label{eq:sz}
    \col([A_1; \cdots ; A_{i+1}]) = \col([A_1 ; A_2 ; \cdots ; A_i ; (A_{i+1} U)|_{[s_{i+1}]}]).
\end{equation}
%
%
Then \eqref{eq:x} and \eqref{eq:sz} together give
    \[
    \col([A_1 ; A_2 ; \cdots ; A_i; A_{i+1}]) = \col([(A_1U)|_{[s_1]} ; \cdots ; (A_iU)|_{[s_i]}; (A_{i+1}U)|_{[s_{i+1}]}]).
    \]
This proves the claim by induction.
\end{proof}



\begin{proof}[Proof of \autoref{th:uncertainty-principle-n}]
We will prove $n_1 \leq \subrank_2(T) \subrank_3(T)$.
Let $A_1, \ldots, A_{n_3} \in \FF^{n_1 \times n_2}$ be the $3$-slices of $T$.
Let~$U$ be as in \autoref{lem:max-rank-induction}.
Let 
\[
s_i = \rank([A_1; \cdots ; A_{i}]) - \rank([A_1; \cdots ; A_{i-1}]).
\]
Conciseness of~$T$ gives that $s_1 + \cdots + s_{n_3} = n_1$.
It follows from this and \autoref{lem:max-rank-induction} that the columns of the submatrices $(A_iU)|_{[s_i]}$ are linearly independent.
This implies that for each $i\in[n_3]$, 
there is a 3-slice of rank at least $s_i$,
%
and in particular $\subrank_3(T)\geq\max_i s_i$. %
%
Let $S$ be the tensor with 3-slices $A_1U, A_2U, \ldots, A_{n_3}U$. Then $\subrank_2(S) \leq \subrank_2(T)$ since the 2-slices of $S$ are linear combinations of the 2-slices of $T$. 
The matrix $M = [(A_1U)|_{[1]};\cdots ; (A_{n_3}U)|_{[1]}]$ consisting of every first column of the 3-slices of $S$ equals the first 2-slice of~$S$. Since the linearly independent columns in the 3-slices of~$S$ are flushed to the left, we find that $M$ has at least $|\{i\in [n_3]: s_i\ne 0\}|$ linearly independent columns, and thus $\rank(M) \geq |\{i\in [n_3]: s_i\ne 0\}|$,
%
%
which gives a lower bound on~$\subrank_2(S)$ and thus on $\subrank_2(T)$.
Since
\[
n_1  = s_1 + \cdots + s_{n_3}\leq |\{i\in [n_3]: s_i\ne 0\}|\, \max_i s_i \leq \subrank_2(T) \subrank_3(T),
\]
the result follows. %\ILnote{I think we should add: "for $\ell_1=2,\ell_2=3,\ell_3=1$ and similarly for other choices." and add to the last equation above the inequality $\leq\subrank_2(T)\subrank_3(T)$}.
%
%
%
%
%
%
%
%
%
%
%
\end{proof}
% Exact native author source LF lines 1731--1731
\setcounter{subsection}{3}\setcounter{theorem}{7}
\subsection{Lower bound on the asymptotic subrank via max-ranks and matrix multiplication tensors}\label{subsec:lb-mamu}
% Exact native author source LF lines 1737--1834

For $a,b,c \in \NN$ we denote by $\langle a,b,c\rangle \in \FF^{ab} \otimes \FF^{bc} \otimes \FF^{ca}$ the %
matrix multiplication tensor $\langle a,b,c \rangle = \sum_{i=1}^a \sum_{j=1}^b \sum_{k=1}^c e_{i,j} \otimes e_{j,k} \otimes e_{k,i}$. 
These tensors have the multiplicative property $\langle a,b,c\rangle \otimes \langle x,y,z\rangle = \langle ax, by, cz\rangle$. They are naturally related to the parameters~$\subrank_i$ as follows:

\begin{lemma}\label{lem:Qi-mamu}
Let $T$ be a tensor and $r \in \NN$. The following holds:
\begin{itemize}
\item $\subrank_1(T) \geq r$ if and only if $T \geq \langle 1, 1,r \rangle$.
\item $\subrank_2(T) \geq r$ if and only if $T \geq \langle r, 1,1 \rangle$.
\item $\subrank_3(T) \geq r$ if and only if $T \geq \langle 1, r, 1 \rangle$.  
\end{itemize}
\end{lemma}

\begin{lemma}\label{lem:Qi-square}
    If there are distinct ${\ell_1},{\ell_2}$ such that $\subrank_{\ell_1}(T), \subrank_{\ell_2}(T) \geq r$, then $\subrank(T^{\boxtimes 2}) \geq r$.
\end{lemma}

\begin{proof}
Suppose $\subrank_1(T), \subrank_2(T) \geq r$ (the proof for the other cases goes the same). Then we have that
$
T \geq \langle1, 1, r\rangle =  \sum_{i=1}^r e_1 \otimes e_i \otimes e_i$ and 
$T \geq \langle r, 1, 1\rangle = \sum_{i=1}^r e_i \otimes e_1 \otimes e_i$.
%
We multiply these inequalities to get
$T^{\boxtimes 2} \geq \langle r,1,r\rangle= \sum_{i,j=1}^r e_i \otimes e_j \otimes (e_i \otimes e_j).
$
We apply the projection $(e_i \otimes e_j) \mapsto \delta_{i=j} e_i$ to the third tensor leg to get
$
T^{\boxtimes 2} \geq \sum_{i=1}^r e_i \otimes e_i \otimes e_i
$
which proves the claim.
\end{proof}

We are now ready to prove the fourth-root lower bound on the asymptotic subrank. After that we will prove the stronger cube-root lower bound (\autoref{th:cube-root-lb}).

\begin{theorem}\label{th:fourth-root}
    Let $T \in \FF^{n_1} \otimes \FF^{n_2} \otimes \FF^{n_3}$ be concise. Then $\asympsubrank(T) \geq (\min_\ell n_\ell)^{1/4}$.
\end{theorem}
\begin{proof}
The proof is an application of \autoref{th:uncertainty-principle-n}.
That theorem only holds over large enough fields however. 
Since the asymptotic subrank does not change under field extension \cite[Theorem~3.10]{strassen1988asymptotic}, we may assume that our base field $\FF$ is infinite (by working over the algebraic closure of the original field) so that the field size condition of \autoref{th:uncertainty-principle-n} is satisfied. %
We use that there are distinct $\ell_1, \ell_2 \in [3]$ such that 
%\refnote{13}{missing square root on $n_i$}\ILnote{maybe it was missing in an earlier version?}\Jshort{yes we fixed this earlier} 
$\subrank_{\ell_1}(T), \subrank_{\ell_2}(T) \geq \sqrt{\min_\ell n_\ell}$ (\autoref{cor:twoLargeRankDirs}).
Then $\asympsubrank(T^{\boxtimes 2}) \geq \sqrt{\min_\ell n_\ell}$ (\autoref{lem:Qi-square}) which gives the claim. %
\end{proof}




%

It turns out that the above construction is not optimal. We will now give a slightly more elaborate construction that leads to the better cube-root bound:

\begin{theorem}\label{th:cube-root-lb}
    Let $T \in \FF^{n_1} \otimes \FF^{n_2} \otimes \FF^{n_3}$ be concise. Then $\asympsubrank(T) \geq (\min_{\ell} n_\ell)^{1/3}$.
\end{theorem}

The proof of \autoref{th:cube-root-lb} makes use of basic properties of matrix multiplication tensors which we discuss in the following two lemmas. 
%\refnote{14}{The following is a repetition of matrix multiplication notation description -- it is used above in this page already} %Recall that we use the notation~$\langle a,b,c\rangle$ for~$a,b,c \in \NN$ to denote the matrix multiplication tensors.
%\Jshort{Removed it.}

\begin{lemma}\label{lem:cube-mamu}
    $T^{\boxtimes 3} \geq \langle \subrank_2(T),\subrank_3(T),\subrank_1(T) \rangle.$
\end{lemma}

\begin{proof}
    From \autoref{lem:Qi-mamu} we have %three inequalities
    %
    %
    %
    %
    %
    %
    $T \geq \langle \subrank_2(T),1,1\rangle$, 
    $T \geq \langle 1,\subrank_3(T),1\rangle$, and  
    $T \geq \langle 1,1,\subrank_1(T)\rangle$.
    %
We multiply these inequalities using the Kronecker product to get the claim.
\end{proof}


\begin{lemma}\label{lem:asympsub-min}
    $\asympsubrank(T^{\boxtimes 3}) \geq \min_{{\ell_1},{\ell_2}} \subrank_{\ell_1}(T) \subrank_{\ell_2}(T).$
\end{lemma}
\begin{proof}
    It is known that for every $a,b,c \in \NN$ the asymptotic subrank of the matrix multiplication tensor $\langle a,b,c\rangle$ is given by $\asympsubrank(\langle a,b,c\rangle) = \min \{ab, bc, ac\}$ \cite[Equation~4.6]{strassen1988asymptotic}. We combine this with \autoref{lem:cube-mamu} to get the claim.
\end{proof}

\begin{proof}[Proof of \autoref{th:cube-root-lb}]
As in the proof of \autoref{th:fourth-root} we may assume that the base field $\FF$ is infinite.
%
%
%
We use $\asympsubrank(T^{\boxtimes 3}) \geq \min_{{\ell_1},{\ell_2}} \subrank_{\ell_1}(T) \subrank_{\ell_2}(T)$ (\autoref{lem:asympsub-min}) and we use that for every distinct ${\ell_1}, {\ell_2}, {\ell_3} \in [3]$ we have $\subrank_{\ell_1}(T)\subrank_{\ell_2}(T) \geq n_{\ell_3}$ (\autoref{th:uncertainty-principle-n}) to obtain
$\asympsubrank(T^{\boxtimes 3}) \geq \min_{{\ell_1},{\ell_2}} \subrank_{\ell_1}(T) \subrank_{\ell_2}(T) \geq \min_{\ell_3} n_{\ell_3}$. %
\end{proof}
\begin{thebibliography}{Str88}
\bibitem[BCS97]{burgisser1997algebraic}
Peter B\"{u}rgisser, Michael Clausen, and Mohammad~Amin Shokrollahi.
\newblock {\em Algebraic complexity theory}, volume 315 of {\em Grundlehren der mathematischen Wissenschaften}.
\newblock Springer-Verlag, Berlin, 1997.
\newblock \href {https://doi.org/10.1007/978-3-662-03338-8} {\path{doi:10.1007/978-3-662-03338-8}}.

\bibitem[CK11]{MR2776439}
Enrico Carlini and Johannes Kleppe.
\newblock Ranks derived from multilinear maps.
\newblock {\em J. Pure Appl. Algebra}, 215(8):1999--2004, 2011.
\newblock \href {https://doi.org/10.1016/j.jpaa.2010.11.010} {\path{doi:10.1016/j.jpaa.2010.11.010}}.

\bibitem[CVZ23]{MR4495838}
Matthias Christandl, P\'{e}ter Vrana, and Jeroen Zuiddam.
\newblock Universal points in the asymptotic spectrum of tensors.
\newblock {\em J. Amer. Math. Soc.}, 36(1):31--79, 2023.
\newblock \href {https://doi.org/10.1090/jams/996} {\path{doi:10.1090/jams/996}}.

\bibitem[DL78]{DEMILLO1978193}
Richard~A. Demillo and Richard~J. Lipton.
\newblock A probabilistic remark on algebraic program testing.
\newblock {\em Information Processing Letters}, 7(4):193--195, 1978.
\newblock \href {https://doi.org/10.1016/0020-0190(78)90067-4} {\path{doi:10.1016/0020-0190(78)90067-4}}.

\bibitem[Sch80]{schwartz1980fast}
Jacob~T. Schwartz.
\newblock Fast probabilistic algorithms for verification of polynomial identities.
\newblock {\em Journal of the ACM}, 27(4):701--717, 1980.
\newblock \href {https://doi.org/10.1145/322217.322225} {\path{doi:10.1145/322217.322225}}.

\bibitem[Str87]{MR882307}
Volker Strassen.
\newblock Relative bilinear complexity and matrix multiplication.
\newblock {\em J. Reine Angew. Math.}, 375/376:406--443, 1987.
\newblock \href {https://doi.org/10.1515/crll.1987.375-376.406} {\path{doi:10.1515/crll.1987.375-376.406}}.

\bibitem[Str88]{strassen1988asymptotic}
Volker Strassen.
\newblock The asymptotic spectrum of tensors.
\newblock {\em J. Reine Angew. Math.}, 384:102--152, 1988.
\newblock \href {https://doi.org/10.1515/crll.1988.384.102} {\path{doi:10.1515/crll.1988.384.102}}.

\bibitem[Str91]{strassen1991degeneration}
Volker Strassen.
\newblock {Degeneration and complexity of bilinear maps: some asymptotic spectra}.
\newblock {\em J. Reine Angew. Math.}, 413:127–180, 1991.
\newblock \href {https://doi.org/10.1515/crll.1991.413.127} {\path{doi:10.1515/crll.1991.413.127}}.

\bibitem[Zip79]{zippel1979probabilistic}
Richard Zippel.
\newblock {\em Probabilistic algorithms for sparse polynomials}.
\newblock Springer, 1979.
\newblock \href {https://doi.org/10.1007/3-540-09519-5_73} {\path{doi:10.1007/3-540-09519-5_73}}.

\end{thebibliography}\end{document}
